\documentclass[11pt]{article}

\usepackage{amsmath, amssymb, amsthm}
\usepackage[margin=1in]{geometry}

\DeclareMathOperator{\disc}{disc}

\title{Compendium of Combinatorics}
\author{DeepSeek prompted by Riyad Derguini}
\date{}

\begin{document}

\maketitle

\textbf{Notation and Conventions.} Let \(\mathbb{N}\) denote the set of positive integers. For \(n \in \mathbb{N}\), we define \([n] = \{1, 2, \dots, n\}\). The cardinality of a set \(S\) is denoted by \(|S|\). The binomial coefficient is denoted by \(\binom{n}{k}\). The falling factorial is \((n)_k = n(n-1)\cdots(n-k+1)\). The natural logarithm is \(\log\), and \(e\) is the base of the natural logarithm.

\hrulefill

\part{Fundamental Counting Principles}

\section{The Basic Rules}
\begin{itemize}
    \item \textbf{Rule of Sum (Addition Principle):} If a task can be performed in \(m\) ways and another disjoint task in \(n\) ways, then there are \(m + n\) ways to perform exactly one of the tasks. Formally, if \(A \cap B = \emptyset\), \(|A \cup B| = |A| + |B|\).
    \item \textbf{Rule of Product (Multiplication Principle):} If a task can be broken down into sequential stages, with \(n_1\) ways to complete stage 1, \(n_2\) ways for stage 2, etc., then the total number of ways is \(\prod_{i=1}^k n_i\).
    \item \textbf{Inclusion-Exclusion Principle (PIE):} For finite sets \(A_1, \dots, A_n\):
    \[
 \left| \bigcup_{i=1}^n A_i \right| = \sum_{i} |A_i| - \sum_{i < j} |A_i \cap A_j| + \sum_{i < j < k} |A_i \cap A_j \cap A_k| - \dots + (-1)^{n-1} |A_1 \cap \dots \cap A_n| 
\]
    \item \textbf{Complementary Counting:} \(|A| = |U| - |A^c|\), where \(U\) is the universal set.
\end{itemize}

\section{Permutations and Combinations}
\begin{itemize}
    \item \textbf{Permutations (Order Matters):} The number of ways to arrange \(k\) distinct objects chosen from \(n\) distinct objects is:
    \[
 P(n, k) = (n)_k = \frac{n!}{(n-k)!} 
\]
    \item \textbf{Combinations (Order Does Not Matter):} The number of ways to choose a subset of \(k\) objects from \(n\) distinct objects is:
    \[
 \binom{n}{k} = \frac{n!}{k!(n-k)!} 
\]
    \item \textbf{Permutations with Repetition:} The number of strings of length \(k\) from an alphabet of size \(n\) is \(n^k\).
    \item \textbf{Permutations of a Multiset:} The number of distinct arrangements of \(n\) objects where there are \(n_1\) objects of type 1, \(n_2\) of type 2, etc., is the multinomial coefficient:
    \[
 \binom{n}{n_1, n_2, \dots, n_k} = \frac{n!}{n_1! n_2! \dots n_k!} 
\]
\end{itemize}

\hrulefill

\part{Binomial Coefficients and Identities}

\section{Core Identities}
\begin{itemize}
    \item \textbf{Symmetry:} \(\binom{n}{k} = \binom{n}{n-k}\).
    \item \textbf{Pascal's Rule:} \(\binom{n}{k} = \binom{n-1}{k-1} + \binom{n-1}{k}\).
    \item \textbf{Absorption Identity:} \(k\binom{n}{k} = n\binom{n-1}{k-1}\).
    \item \textbf{Vandermonde's Identity:} \(\binom{m+n}{r} = \sum_{k=0}^r \binom{m}{k}\binom{n}{r-k}\).
    \item \textbf{Hockey-Stick Identity:} \(\sum_{i=r}^n \binom{i}{r} = \binom{n+1}{r+1}\).
\end{itemize}

\section{The Binomial Theorem}
For any \(x, y \in \mathbb{R}\) and \(n \in \mathbb{N}\):
\[
 (x + y)^n = \sum_{k=0}^n \binom{n}{k} x^k y^{n-k} 
\]
\begin{itemize}
    \item \textbf{Corollary 1:} \(\sum_{k=0}^n \binom{n}{k} = 2^n\).
    \item \textbf{Corollary 2:} \(\sum_{k=0}^n (-1)^k \binom{n}{k} = 0\) for \(n \ge 1\).
\end{itemize}

\section{The Multinomial Theorem}
\[
 (x_1 + x_2 + \dots + x_m)^n = \sum_{k_1 + \dots + k_m = n} \binom{n}{k_1, \dots, k_m} x_1^{k_1} \dots x_m^{k_m} 
\]

\hrulefill

\part{Stars and Bars (Distributions)}

The number of ways to distribute \(n\) identical items into \(k\) distinct bins:
\begin{enumerate}
    \item \textbf{No restrictions (non-negative solutions):}
    \[
 \binom{n + k - 1}{k - 1} 
\]
    \item \textbf{Each bin gets at least one item (positive solutions):}
    \[
 \binom{n - 1}{k - 1} 
\]
    \item \textbf{Each bin gets at least \(r\) items:}
    \[
 \binom{n - rk + k - 1}{k - 1} 
\]
\end{enumerate}

\hrulefill

\part{Recurrence Relations \& Generating Functions}

\section{CommonRecurrences}
\begin{itemize}
    \item \textbf{Fibonacci Sequence:} \(F_n = F_{n-1} + F_{n-2}\), with \(F_0=0, F_1=1\). Closed form (Binet's Formula): \(F_n = \frac{1}{\sqrt{5}}\left(\left(\frac{1+\sqrt{5}}{2}\right)^n - \left(\frac{1-\sqrt{5}}{2}\right)^n\right)\).
    \item \textbf{Catalan Numbers:} \(C_n = \frac{1}{n+1}\binom{2n}{n}\). They satisfy \(C_n = \sum_{i=0}^{n-1} C_i C_{n-1-i}\). They count binary trees, valid parenthesis sequences, and polygon triangulations.
    \item \textbf{Derangements (\(D_n\)):} Permutations with no fixed points. \(D_n = n! \sum_{k=0}^n \frac{(-1)^k}{k!} \approx \frac{n!}{e}\).
\end{itemize}

\section{Generating Functions}
\begin{itemize}
    \item \textbf{Ordinary Generating Function (OGF):} \(G(x) = \sum_{n=0}^\infty a_n x^n\).
    \begin{itemize}
        \item Used for sequences where order doesn't matter (combinations).
        \item \emph{Example:} For \(a_n = 1\), \(G(x) = \frac{1}{1-x}\).
    \end{itemize}
    \item \textbf{Exponential Generating Function (EGF):} \(E(x) = \sum_{n=0}^\infty a_n \frac{x^n}{n!}\).
    \begin{itemize}
        \item Used for sequences where order matters (permutations).
        \item \emph{Example:} For \(a_n = n!\), \(E(x) = \frac{1}{1-x}\).
    \end{itemize}
\end{itemize}

\hrulefill

\part{Pigeonhole Principle \& Extremal Combinatorics}

\section{Pigeonhole Principle}
\begin{itemize}
    \item \textbf{Basic:} If \(n\) items are placed into \(m\) containers and \(n > m\), at least one container must contain at least \(\lceil n/m \rceil\) items.
    \item \textbf{Generalized:} If \(N\) objects are placed into \(k\) boxes, at least one box contains \(\lceil N/k \rceil\) objects.
    \item \textbf{Erdős–Szekeres Theorem:} Any sequence of \(n^2 + 1\) distinct real numbers contains a monotonic subsequence of length \(n+1\).
\end{itemize}

\section{Ramsey Theory}
\begin{itemize}
    \item \textbf{Ramsey's Theorem:} For any \(r, s \ge 1\), there exists a minimum integer \(R(r,s)\) such that any 2-coloring of the edges of the complete graph \(K_{R(r,s)}\) contains a red \(K_r\) or a blue \(K_s\).
    \item \textbf{Known Values:} \(R(3,3) = 6\), \(R(3,4) = 9\), \(R(4,4) = 18\).
\end{itemize}

\section{Extremal Graph Theory}
\begin{itemize}
    \item \textbf{Turán's Theorem:} The maximum number of edges in an \(n\)-vertex graph that does not contain a clique of size \(r+1\) is given by the Turán graph \(T(n,r)\), with exactly \(\left(1 - \frac{1}{r}\right)\frac{n^2}{2} + O(n)\) edges.
\end{itemize}

\hrulefill

\part{The Probabilistic Method}

A non-constructive technique to prove the existence of a combinatorial object by showing it has a positive probability of being generated randomly.

\begin{itemize}
    \item \textbf{Union Bound:} \(\mathbb{P}(\bigcup_{i=1}^n A_i) \le \sum_{i=1}^n \mathbb{P}(A_i)\). 
    \begin{itemize}
        \item \emph{Use Case:} If the sum of probabilities of ``bad events'' is strictly less than 1, then there exists an outcome where no bad events occur.
    \end{itemize}
    \item \textbf{Linearity of Expectation:} \(\mathbb{E}[\sum X_i] = \sum \mathbb{E}[X_i]\). (Does not require independence).
    \begin{itemize}
        \item \emph{Use Case:} If \(\mathbb{E}[X] = \mu\), then there must exist an outcome where \(X \ge \mu\) and an outcome where \(X \le \mu\).
    \end{itemize}
    \item \textbf{Lovász Local Lemma (LLL):} Let \(A_1, \dots, A_n\) be ``bad events'' such that each \(A_i\) is mutually independent of all but at most \(d\) other events. If \(\mathbb{P}(A_i) \le p\) for all \(i\) and \(e \cdot p \cdot (d+1) \le 1\) (where \(e\) is the base of the natural log), then \(\mathbb{P}(\bigcap \overline{A_i}) > 0\).
\end{itemize}

\hrulefill

\part{Discrepancy Theory}

\begin{itemize}
    \item \textbf{Standard Definition:} For a set system \(\mathcal{F}\) on a ground set \(V\), the discrepancy is \(\disc(\mathcal{F}) = \min_{\chi: V \to \{-1,1\}} \max_{S \in \mathcal{F}} |\sum_{x \in S} \chi(x)|\).
    \item \textbf{Matrix Form:} Minimizing \(\|A\mathbf{x}\|_\infty\) where \(A\) is the incidence matrix and \(\mathbf{x} \in \{\pm 1\}^n\).
    \item \textbf{Spencer's Theorem (Six Standard Deviations):} For any set system with \(n\) elements and \(m\) sets, there exists a coloring with discrepancy at most \(O(\sqrt{m \log(n/m)})\). If \(m=n\), the discrepancy is \(O(\sqrt{n})\).
\end{itemize}

\hrulefill

\part{Asymptotics \& Stirling's Approximation}

\begin{itemize}
    \item \textbf{Stirling's Approximation:}
    \[
 n! \sim \sqrt{2\pi n} \left(\frac{n}{e}\right)^n 
\]
    Equivalently:
    \[
 \log(n!) = n \log n - n + O(\log n) 
\]
    \item \textbf{Approximating Binomial Coefficients:}
    \[
 \binom{n}{k} \approx \frac{n^n}{k^k (n-k)^{n-k}} \sqrt{\frac{n}{2\pi k(n-k)}} 
\]
    \item \textbf{Entropy Bound:} For \(0 < p < 1\), where \(k = pn\):
    \[
 \binom{n}{pn} \approx 2^{n H(p)} 
\]
    where \(H(p) = -p \log_2 p - (1-p) \log_2 (1-p)\) is the binary entropy function.
\end{itemize}

\end{document} 
